%!TEX TS-program = lualatex
%-------------------------------------------------------------------------------
%  Curriculum Vitae — Nils Grünefeld
%  Built on Awesome-CV (https://github.com/posquit0/Awesome-CV), class file
%  vendored as awesome-cv.cls. LuaLaTeX only.
%
%  Build:   cd cv && latexmk        (latexmkrc selects LuaLaTeX; writes cv.pdf)
%-------------------------------------------------------------------------------
\documentclass[11pt, a4paper]{awesome-cv}

% Page geometry
\geometry{left=1.6cm, top=1.0cm, right=1.6cm, bottom=1.8cm, footskip=.6cm}

% One restrained accent: the website's navy (--accent on n.grnfld.me)
\definecolor{awesome-navy}{HTML}{1F4E79}
\colorlet{awesome}{awesome-navy}
\setbool{acvSectionColorHighlight}{false}   % plain section titles
\usepackage{needspace}                      % keep section headings with their first entry
\let\acvOrigSection\cvsection
\renewcommand{\cvsection}[1]{\Needspace*{4cm}\acvOrigSection{#1}}

\renewcommand{\acvHeaderSocialSep}{\quad\textbar\quad}

% Awesome-CV's one-line tables (honors, talks, schools) give the right-hand
% location column only 2.5cm; widen it so "City, Country" fits on one line.
\renewenvironment{cvhonors}{%
  \vspace{\acvSectionContentTopSkip}
  \vspace{-2mm}
  \begin{center}
    \setlength\tabcolsep{0pt}
    \setlength{\extrarowheight}{0pt}
    \begin{tabular*}{\textwidth}{@{\extracolsep{\fill}} C{1.5cm} L{\textwidth - 4.8cm} R{3.3cm}}
}{%
    \end{tabular*}
  \end{center}
}

%-------------------------------------------------------------------------------
%  Personal information
%-------------------------------------------------------------------------------
\name{Nils}{Grünefeld}
\position{PhD Fellow{\enskip\cdotp\enskip}Uncertainty Quantification in Machine Learning}
\address{NLPnorth, IT University of Copenhagen{\enskip\cdotp\enskip}Copenhagen, Denmark}
\mobile{(+45) 50 12 12 82}
\email{nilgr@itu.dk}
\homepage{n.grnfld.me}
\github{ngruenefeld}
\linkedin{ngruenefeld}
\googlescholar{sG390YIAAAAJ}{}

\begin{document}
\makecvheader[C]
\makecvfooter
  {\today}
  {Nils~Grünefeld~~~·~~~Curriculum Vitae}
  {\thepage}

%===============================================================================
\cvsection{Education}
\begin{cventries}
  \cventry
    {PhD in Data Science}
    {IT University of Copenhagen}
    {Copenhagen, Denmark}
    {Oct. 2025 -- Present}
    {
      \begin{cvitems}
        \item {Thesis: \emph{Efficient Gradient-Based Uncertainty Quantification using Language Perturbation}.}
        \item {Supervised by Christian~Hardmeier and Jes~Frellsen. NLPnorth group; affiliated with the Pioneer Centre for AI.}
        \item {Funded by a Danish Data Science Academy (DDSA) PhD Fellowship.}
      \end{cvitems}
    }
  \cventry
    {MSc in Computational Linguistics}
    {Ludwig Maximilian University of Munich}
    {Munich, Germany}
    {2023 -- 2025}
    {
      \begin{cvitems}
        \item {Thesis: \emph{Estimating Uncertainty in Natural Language Processing}, supervised by Barbara~Plank and Silvia~Casola (MaiNLP).}
      \end{cvitems}
    }
  \cventry
    {MSc in Computer Science}
    {Leibniz University Hannover}
    {Hannover, Germany}
    {2022 -- 2023}
    {
      \begin{cvitems}
        \item {Transferred to the MSc in Computational Linguistics at LMU Munich after one semester.}
      \end{cvitems}
    }
  \cventry
    {BSc in Computer Science}
    {Leibniz University Hannover}
    {Hannover, Germany}
    {2014 -- 2022}
    {
      \begin{cvitems}
        \item {Thesis: \emph{Explaining and Applying Graph Neural Networks on Text}, supervised by Avishek~Anand and Thorben~Funke.}
      \end{cvitems}
    }
\end{cventries}

%===============================================================================
\cvsection{Publications}
\par
\begin{cvparagraph}
  $^{*}$\,Equal contribution.
\end{cvparagraph}
\cvsubsection{Peer-Reviewed}
\begin{cventries}
  \cventry
    {\textbf{Nils~Grünefeld}, Jes~Frellsen, Christian~Hardmeier}
    {\href{https://proceedings.mlr.press/v327/grunefeld26a.html}{An Isotropic Approach to Efficient Uncertainty Quantification with Gradient Norms}}
    {ProbML 2026{\enskip\cdotp\enskip}Best Paper Award}
    {2026}
    {
      \begin{cvitems}
        \item {Proceedings of the 1st Symposium on Probabilistic Machine Learning, PMLR 327:86--118. Oral presentation; also presented as a poster at the FoGen Workshop, ICML 2026.}
      \end{cvitems}
    }
\end{cventries}
\cvsubsection{Under Review}
\begin{cventries}
  \cventry
    {Johannes~Gabriel~Sindlinger$^{*}$, \textbf{Nils~Grünefeld}$^{*}$, Anna~Rogers, Jes~Frellsen, Christian~Hardmeier}
    {Separating Ambiguity from Ignorance: Explaining Data Influence by Uncertainty~Reduction}
    {Under review{\enskip\cdotp\enskip}AISTATS 2027}
    {2026}
    {}
  \vspace{2.5mm}
  \cventry
    {\textbf{Nils~Grünefeld}$^{*}$, Bertram~Højer$^{*}$, Philipp~Mondorf, Barbara~Plank, Anna~Rogers, Christian~Hardmeier, Stefan~Heinrich, Jes~Frellsen}
    {\href{https://arxiv.org/abs/2605.07776}{Tracing Uncertainty in Language Model ``Reasoning''}}
    {Under review{\enskip\cdotp\enskip}ICLR 2027}
    {2026}
    {
      \begin{cvitems}
        \item {arXiv:2605.07776. Presented as a poster at the FoGen Workshop, ICML 2026.}
      \end{cvitems}
    }
\end{cventries}
\cvsubsection{In Preparation}
\begin{cventries}
  \cventry
    {\textbf{Nils~Grünefeld}, Mattes~Ruckdeschel, Silvia~Casola, Marco~Stranisci, Anna~Rogers, Barbara~Plank, Jes~Frellsen, Christian~Hardmeier}
    {How do LLMs learn and express different perspectives?}
    {}
    {2026}
    {}
\end{cventries}

%===============================================================================
\cvsection{Honors \& Awards}
\begin{cvhonors}
  \cvhonor
    {Best Paper Award}
    {ProbML 2026}
    {Seoul, South Korea}
    {2026}
  \cvhonor
    {PhD Fellowship}
    {Danish Data Science Academy (DDSA)}
    {Denmark}
    {2025}
\end{cvhonors}

%===============================================================================
\cvsection{Talks}
\begin{cventries}
  \cventry
    {Invited Talk{\enskip\cdotp\enskip}Mathematical Data Science group, University of Vienna}
    {Separating Ambiguity from Ignorance: Explaining Data Influence by Uncertainty~Reduction}
    {Vienna, Austria}
    {Nov. 2026}
    {}
  \vspace{2.5mm}
  \cventry
    {Invited Talk{\enskip\cdotp\enskip}Andrés~Masegosa's group, Aalborg University Copenhagen}
    {An Isotropic Approach to Efficient Uncertainty Quantification with Gradient Norms}
    {Copenhagen, Denmark}
    {Sep. 2026}
    {}
  \vspace{2.5mm}
  \cventry
    {Oral{\enskip\cdotp\enskip}ProbML 2026}
    {An Isotropic Approach to Efficient Uncertainty Quantification with Gradient Norms}
    {Seoul, South Korea}
    {Jul. 2026}
    {}
\end{cventries}

%===============================================================================
\cvsection{Experience}
\begin{cventries}
  \cventry
    {Data Scientist}
    {the peak lab.}
    {Oldenburg, Germany}
    {2022 -- 2025}
    {
      \begin{cvitems}
        \item {Analyzed complex data sets using statistical methods and machine learning to extract insights, inform decisions, and develop data-driven products and solutions.}
        \item {Designed and implemented a standardized process for data science projects.}
        \item {Implemented a company-wide focus on explainability in machine learning solutions.}
        \item {Designed and implemented machine learning solutions for major company clients, including text classification for \mbox{business product categorization} and time-series-based predictive maintenance modeling.}
      \end{cvitems}
    }
\end{cventries}

%===============================================================================
\cvsection{Teaching \& Supervision}
\begin{cventries}
  \cventry
    {Guest Lecture: Uncertainty in LLM Reasoning}
    {Reasoning in Large Language Models (MSc seminar)}
    {LMU Munich}
    {Autumn 2026}
    {}
  \vspace{2.5mm}
  \cventry
    {Lecturer}
    {Advanced Natural Language Processing (MSc course)}
    {IT University of Copenhagen}
    {Autumn 2026}
    {}
  \vspace{2.5mm}
  \cventry
    {Co-Supervisor (with Christian~Hardmeier)}
    {Uncertainty Communication in LLMs: Can Fine-tuning Improve Calibration in Question~Answering? (BSc thesis)}
    {IT University of Copenhagen}
    {Spring 2026}
    {}
  \vspace{2.5mm}
  \cventry
    {Teaching Assistant}
    {Programming I \& II (BSc courses)}
    {Leibniz University Hannover}
    {2016 -- 2021}
    {}
\end{cventries}

%===============================================================================
\cvsection{Academic Service}
\begin{cvskills}
  \cvskill
    {Reviewing}
    {Computational Linguistics (MIT Press){\enskip\cdotp\enskip}Transactions on Machine Learning Research (TMLR)}
\end{cvskills}

%===============================================================================
\cvsection{Summer \& Winter Schools}
\begin{cvhonors}
  \cvhonor
    {Machine Learning Summer School (MLSS)}
    {}
    {New York City, USA}
    {2026}
  \cvhonor
    {Circle U, NLPL \& OpenEuroLLM Winter School on Multilinguality in LLMs}
    {}
    {Skeikampen, Norway}
    {2026}
  \cvhonor
    {Athens Natural Language Processing Summer School}
    {}
    {Athens, Greece}
    {2024}
\end{cvhonors}

%===============================================================================
\cvsection{Skills}
\begin{cvskills}
  \cvskill
    {Programming}
    {Python, Java, C, R, Bash, Git, \LaTeX}
  \cvskill
    {Frameworks}
    {PyTorch, Hugging Face Transformers, Pyro, (Py)Stan, scikit-learn, NumPy, SciPy, pandas, Matplotlib}
  \cvskill
    {Languages}
    {German (native), English (fluent), French (basic)}
\end{cvskills}

\end{document}
